\chapter{2014年广东卷（文）}

\section{单选题}

\begin{problem}
已知集合 \(M = \{2, 3, 4\}\)，\(N = \{0, 2, 3, 5\}\)，则 \(M \cap N\)（\quad）

\choices
    {\(\{0, 2\}\)}
    {\(\{2, 3\}\)}
    {\(\{3, 4\}\)}
    {\(\{3, 5\}\)}
\end{problem}

\begin{answer}
B
\end{answer}

\begin{problem}
已知复数 \(z\) 满足 \((3 - 4\i)z = 25\)，则 \(z =\)（\quad）

\choices
    {\(-3 - 4\i\)}
    {\(-3 + 4\i\)}
    {\(3 - 4\i\)}
    {\(3 + 4\i\)}
\end{problem}

\begin{answer}
D
\end{answer}

\begin{problem}
已知向量 \(\bs{a} = (1, 2)\)，\(\bs{b} = (3, 1)\)，则 \(\bs{b} - \bs{a} =\)（\quad）

\choices
    {\((-2, 1)\)}
    {\((2, -1)\)}
    {\((2, 0)\)}
    {\((4, 3)\)}
\end{problem}

\begin{answer}
B
\end{answer}

\begin{problem}
若变量 \(x\)，\(y\) 满足约束条件 \(\begin{cases}
x + 2y \leqslant 8,\\
0 \leqslant x \leqslant 4,\\
0 \leqslant y \leqslant 3,
\end{cases}\) 则 \(z = 2x + y\) 的最大值等于（\quad）

\choices
    {7}
    {8}
    {10}
    {11}
\end{problem}

\begin{answer}
C
\end{answer}

\begin{problem}
下列函数为奇函数的是（\quad）

\choices
    {\(2^{x} - \frac{1}{2^{x}}\)}
    {\(x^{3}\sin x\)}
    {\(2\cos x + 1\)}
    {\(x^{2} + 2^{x}\)}
\end{problem}

\begin{answer}
A
\end{answer}

\begin{problem}
为了解 1000 名学生的学习情况，采用系统抽样的方法，从中抽取容量为 40 的样本，则分段的间隔为（\quad）

\choices
    {50}
    {40}
    {25}
    {20}
\end{problem}

\begin{answer}
C
\end{answer}

\begin{problem}
在 \(\triangle ABC\) 中，角 \(A\)，\(B\)，\(C\) 所对应的边分别为 \(a\)，\(b\)，\(c\)，则 “\(a \leqslant b\)” 是 “\(\sin A \leqslant \sin B\)” 的（\quad）

\choices
    {充分必要条件}
    {充分非必要条件}
    {必要非充分条件}
    {非充分非必要条件}
\end{problem}

\begin{answer}
A
\end{answer}

\begin{problem}
若实数 \(k\) 满足 \(0 < k < 5\)，则曲线 \(\frac{x^{2}}{16} - \frac{y^{2}}{5-k} = 1\) 与曲线 \(\frac{x^{2}}{16-k} - \frac{y^{2}}{5} = 1\) 的（\quad）

\choices
    {实半轴长相等}
    {虚半轴长相等}
    {离心率相等}
    {焦距相等}
\end{problem}

\begin{answer}
D
\end{answer}

\begin{problem}
若空间中四条两两不同的直线 \(l_{1}\)，\(l_{2}\)，\(l_{3}\)，\(l_{4}\)，满足 \(l_{1} \perp l_{2}\)，\(l_{2} \parallel l_{3}\)，\(l_{3} \perp l_{4}\)，则下列结论一定正确的是（\quad）

\choices
    {\(l_{1} \perp l_{4}\)}
    {\(l_{1} \parallel l_{4}\)}
    {\(l_{1}\)，\(l_{4}\) 既不垂直也不平行}
    {\(l_{1}\)，\(l_{4}\) 的位置关系不确定}
\end{problem}

\begin{answer}
D
\end{answer}

\begin{problem}
对任意复数 \(w_{1}\)，\(w_{2}\)，定义 \(w_{1} * w_{2} = w_{1} \overline{w}_{2}\). 其中 \(\overline{w}_{2}\) 是 \(w_{2}\) 的共轭复数. 对任意复数 \(z_{1}\)，\(z_{2}\)，\(z_{3}\) 有如下四个命题：

\begin{circlelist}
\item \((z_{1} + z_{2}) * z_{3} = (z_{1} * z_{3}) + (z_{2} * z_{3})\)；
\item \(z_{1} * (z_{2} + z_{3}) = (z_{1} * z_{2}) + (z_{1} * z_{3})\)；
\item \((z_{1} * z_{2}) * z_{3} = z_{1} * (z_{2} * z_{3})\)；
\item \(z_{1} * z_{2} = z_{2} * z_{1}\).
\end{circlelist}

则真命题的个数是（\quad）

\choices
    {1}
    {2}
    {3}
    {4}
\end{problem}

\begin{answer}
B
\end{answer}

\section{填空题}

\begin{problem}
曲线 \(y = -5\e^{x} + 3\) 在点 \((0, -2)\) 处的切线方程为\fillinblank{}.
\end{problem}

\begin{answer}
\(5x + y + 2 = 0\)
\end{answer}

\begin{problem}
从字母 \(a\)，\(b\)，\(c\)，\(d\)，\(e\) 中任取两个不同字母，则取到字母 \(a\) 的概率为\fillinblank{}.
\end{problem}

\begin{answer}
\(\dfrac{2}{5}\)
\end{answer}

\begin{problem}
等比数列 \(\{a_{n}\}\) 的各项均为正数，且 \(a_{1}a_{5} = 4\)，则 \(\log_{2} a_{1} + \log_{2} a_{2} + \log_{2} a_{3} + \log_{2} a_{4} + \log_{2} a_{5} =\fillinblank{}.\)
\end{problem}

\begin{answer}
5
\end{answer}

\section{选考题：填空题}

\begin{problem}
在极坐标系中，曲线 \(C_{1}\) 和 \(C_{2}\) 的方程分别为 \(2\rho \cos^{2}\theta = \sin\theta\) 和 \(\rho \cos\theta = 1\)，以极点为平面直角坐标系的原点，极轴为 \(x\) 轴的正半轴，建立平面直角坐标系，则曲线 \(C_{1}\) 和 \(C_{2}\) 的交点的直角坐标为\fillinblank{}.
\end{problem}

\begin{answer}
\((1, 2)\)
\end{answer}

\begin{problem}
如图，在平行四边形 \(ABCD\) 中，点 \(E\) 在 \(AB\) 上，且 \(EB = 2AE\)，\(AC\) 与 \(DE\) 交于点 \(F\)，则 \(\triangle CDF\) 的周长与 \(\triangle AEF\) 的周长之比为\fillinblank{}.

\bitmapfigure[max width=0.34\linewidth]{img/2014/guangdong_liberal/q15_fig1.png}
\end{problem}

\begin{answer}
3
\end{answer}

\section{解答题}

\begin{problem}
已知函数 \(f(x) = A \sin \left( x + \frac{\pi}{3} \right)\)，\(x \in \R\)，且 \(f\left( \frac{5\pi}{12} \right) = \frac{3\sqrt{2}}{2}\).

\begin{enumerate}
\item 求 \(A\) 的值；
\item 若 \(f(\theta) - f(-\theta) = \sqrt{3}\)，\(\theta \in \left(0, \frac{\pi}{2}\right)\)，求 \(f\left(\frac{\pi}{6} - \theta\right)\).
\end{enumerate}
\end{problem}

\begin{solution}
\begin{enumerate}
\item 短答案：\(A=3\).
\item 短答案：\(f\left(\frac{\pi}{6}-\theta\right)=\sqrt{6}\).
\end{enumerate}
\end{solution}

\begin{problem}
某车间 20 名工人年龄数据如下表：

\begin{center}
\begin{tabular}{c|c}
\hline
年龄（岁） & 工人数（人） \\
\hline
19 & 1 \\
28 & 3 \\
29 & 3 \\
30 & 5 \\
31 & 4 \\
32 & 3 \\
40 & 1 \\
\hline
合计 & 20 \\
\hline
\end{tabular}
\end{center}

\begin{enumerate}
\item 求这 20 名工人年龄的众数与极差；
\item 以十位数为茎，个位数为叶，作出这 20 名工人年龄的茎叶图；
\item 求这 20 名工人年龄的方差.
\end{enumerate}
\end{problem}

\begin{solution}
\begin{enumerate}
\item 短答案：众数为 \(30\)，极差为 \(21\).
\item 短答案：茎叶图为
\begin{center}
\begin{tabular}{c|l}
茎 & 叶 \\
\hline
1 & 9 \\
2 & \(8\quad 8\quad 8\quad 9\quad 9\quad 9\) \\
3 & \(0\quad 0\quad 0\quad 0\quad 0\quad 1\quad 1\quad 1\quad 1\quad 2\quad 2\quad 2\) \\
4 & 0
\end{tabular}
\end{center}
\item 短答案：方差为 \(12.6\).
\end{enumerate}
\end{solution}

\begin{problem}
如图 1，四边形 \(ABCD\) 为矩形，\(PD \perp\) 平面 \(ABCD\)，\(AB = 1\)，\(BC = PC = 2\)，作如图 2 折叠，折痕 \(EF \parallel DC\). 其中点 \(E\)，\(F\) 分别在线段 \(PD\)，\(PC\) 上，沿 \(EF\) 折叠后点 \(P\) 叠在线段 \(AD\) 上的点记为 \(M\)，并且 \(MF \perp CF\).

\begin{enumerate}
\item 证明：\(CF \perp\) 平面 \(MDF\)；
\item 求三棱锥 \(M-CDE\) 的体积.
\end{enumerate}

\examfiguregroup[float=inline]{img/2014/guangdong_liberal/q18_fig1.png}{%
    \makebox[0.44\linewidth][c]{\bitmapinclude[max width=\linewidth]{img/2014/guangdong_liberal/q18_fig1.png}}\hfill
    \makebox[0.44\linewidth][c]{\bitmapinclude[max width=\linewidth]{img/2014/guangdong_liberal/q18_fig2.png}}%
}
\end{problem}

\begin{solution}
\begin{enumerate}
\item 短答案：\(CF \perp\) 平面 \(MDF\).
\item 短答案：三棱锥 \(M-CDE\) 的体积为 \(\dfrac{1}{3}\).
\end{enumerate}
\end{solution}

\begin{problem}
设各项均为正数的数列 \(\{a_{n}\}\) 的前 \(n\) 项和为 \(S_{n}\)，且 \(S_{n}\) 满足 \(S_{n}^{2} - (n^{2} + n - 3)S_{n} - 3(n^{2} + n) = 0\)，\(n \in \N^{*}\).

\begin{enumerate}
\item 求 \(a_{1}\) 的值；
\item 求数列 \(\{a_{n}\}\) 的通项公式；
\item 证明：对一切正整数 \(n\)，有 \(\frac{1}{a_1(a_1 + 1)} + \frac{1}{a_2(a_2 + 1)} + \cdots + \frac{1}{a_n(a_n + 1)} < \frac{1}{3}\).
\end{enumerate}
\end{problem}

\begin{solution}
\begin{enumerate}
\item 短答案：\(a_{1}=2\).
\item 短答案：\(a_{n}=2n\).
\item 短答案：\(\displaystyle\sum_{k=1}^{n}\frac{1}{a_k(a_k+1)}<\frac{1}{3}\).
\end{enumerate}
\end{solution}

\begin{problem}
已知椭圆 \(C: \frac{x^2}{a^2} + \frac{y^2}{b^2} = 1\)（\(a > b > 0\)）的一个焦点为 \((\sqrt{5}, 0)\)，离心率为 \(\frac{\sqrt{5}}{3}\).

\begin{enumerate}
\item 求椭圆 \(C\) 的标准方程；
\item 若动点 \(P(x_0, y_0)\) 为椭圆 \(C\) 外一点，且点 \(P\) 到椭圆 \(C\) 的两条切线互相垂直，求点 \(P\) 的轨迹方程.
\end{enumerate}
\end{problem}

\begin{solution}
\begin{enumerate}
\item 短答案：椭圆 \(C\) 的标准方程为 \(\dfrac{x^{2}}{9}+\dfrac{y^{2}}{4}=1\).
\item 短答案：点 \(P\) 的轨迹方程为 \(x_{0}^{2}+y_{0}^{2}=13\).
\end{enumerate}
\end{solution}

\begin{problem}
已知函数 \(f(x) = \frac{1}{3}x^3 + x^2 + ax + 1\)（\(a \in \R\)）.

\begin{enumerate}
\item 求函数 \(f(x)\) 的单调区间；
\item 当 \(a < 0\) 时，试讨论是否存在 \(x_0 \in \left(0, \frac{1}{2}\right) \cup \left(\frac{1}{2}, 1\right)\)，使得 \(f(x_0) = f\left(\frac{1}{2}\right)\).
\end{enumerate}
\end{problem}

\begin{solution}
\begin{enumerate}
\item 短答案：当 \(a \geqslant 1\) 时，函数在 \(\R\) 上单调递增；当 \(a < 1\) 时，函数的单调递增区间为 \(\left(-\infty,-1-\sqrt{1-a}\right)\) 和 \(\left(-1+\sqrt{1-a},+\infty\right)\)，单调递减区间为 \(\left(-1-\sqrt{1-a},-1+\sqrt{1-a}\right)\).
\item 短答案：当 \(a \in \left(-\dfrac{25}{12},-\dfrac{5}{4}\right) \cup \left(-\dfrac{5}{4},-\dfrac{7}{6}\right)\) 时存在这样的 \(x_{0}\)，否则不存在.
\end{enumerate}
\end{solution}
